\section{Canonical nominal actions and support}
\label{sec:canonical-support}

The selected actions on atoms, discrete data, products and finite atom sets
give basic contexts for nominal reasoning. Their nominality is a property
of those actions, not an additional choice of action.

\begin{proposition}[Canonical nominality]
\label{prop:canonical-nominality}
For any atom carrier $A$, the canonical action on $A$ and the trivial action
on $\Discrete_A(X)$ are nominal. Products of nominal actions are nominal
under the componentwise action, and $\Finset(A)$ is nominal under direct image.
\end{proposition}

\begin{proof}
The singleton $\{a\}$ supports the atom $a$, since a permutation fixing
that singleton fixes its member. Every permutation fixes a discrete value,
so the empty set supports it. The union of supporting bounds for $x$ and
$y$ supports $(x,y)$ by Proposition~\ref{prop:support-products}.
Finally, a permutation fixing every member of $S$ has image $S$ itself,
so $S$ supports $S\in\Finset(A)$.
\end{proof}

These certificates need no infinitude assumption. In particular, the
discrete construction assumes no action or nominality of the underlying
data type $X$. The Lean proofs reuse Mathlib's support-of-a-member and
finite-image laws, together with the existing finite-supportedness product
law. The Finset certificate retains the equality assumption of its image
action and is registered in the same \texttt{Pointwise} scope. The other
three certificates need no decidable equality. All certify the selected
actions of Section~\ref{sec:actions}; none introduces a competing action.

\begin{proposition}[Exact canonical supports]
\label{prop:canonical-support}
Assume $A$ is infinite. Under these canonical actions,
\begin{align*}
  \operatorname{supp}_A(a)&=\{a\}, &
  \operatorname{supp}_A(\operatorname{mk}(d))&=\varnothing,\\
  \operatorname{supp}_A(x,y)&=
    \operatorname{supp}_A(x)\cup\operatorname{supp}_A(y), &
  \operatorname{supp}_A(S)&=S.
\end{align*}
The product equation requires only finite supportedness of the two
particular elements, rather than nominality of their entire carriers.
\end{proposition}

\begin{proof}
The bounds in Proposition~\ref{prop:canonical-nominality} give the
upper inclusions by leastness. For an atom, suppose a finite supporting
bound $T$ omits $a$. Choose $b\notin T\cup\{a\}$. The swap $\swap a b$
fixes $T$ pointwise but sends $a$ to the distinct atom $b$, a contradiction.
Thus every supporting bound contains $a$. For a discrete value the empty
bound is already least.

A bound supporting $(x,y)$ supports each projection. Consequently, each
component's least support is contained in the pair's least support, giving
the reverse product inclusion. Conversely, their union supports the pair.

For the finite-set equation, suppose $T$ supports $S$ and omits an atom
$a\in S$. Choose $b\notin S\cup T$. Both endpoints of $\swap a b$ lie
outside $T$, so the swap fixes $S$ by the support swap criterion. But its
image contains $b$, the image of $a$, contradicting $b\notin S$. Thus
$S$ is contained in every finite supporting bound and is itself the least.
\end{proof}

The reverse inclusions are essential: a convenient supporting bound need
not identify all atoms on which an element depends. The Lean theorem
\path{FinitelySupported.support_prod} states the elementwise product
equation; carrier equations, including \path{support_finset}, use the
canonical certificates. Agreement of support certificates transfers each
canonical formula to independently supplied evidence of the same element.
Singleton and empty Finset expressions need no decidable equality, whereas
the displayed union and the selected Finset image action retain it.
Classical equality used only inside the atom proof adds no public premise.

Infinitude is a hypothesis of the general least-support interface, distinct
from the preceding nominality certificates. In particular, the canonical
Bool action is nominal while false has no least support, as shown in
Example~\ref{ex:finite-atom-intersection} and
Section~\ref{sec:least-support}.

\begin{example}[An unordered pair is not strongly supported]
\label{ex:unordered-pair-support}
Let $a\ne b$ and $S=\{a,b\}$ over infinite atoms. Then
\begin{gather*}
  \swap a b\act S=S,\qquad \operatorname{supp}_A(S)=\{a,b\},
  \\
  \swap a b(a)=b\ne a,\qquad \swap a b(b)=a\ne b.
\end{gather*}
Thus a permutation can fix an element while moving both atoms of its
least support. The exact Finset formula identifies the actual least
support in this example, not merely a sufficient bound. Fixation of an
element entails setwise preservation of its least support by transport;
it does not entail pointwise fixation.
\end{example}

\section{Freshness and renaming}
\label{sec:freshness}

Assume throughout this section that $A$ is infinite. For finitely supported
elements $x\in X$ and $y\in Y$ of selected actions, define freshness by
\[
  x\mathrel{\#}y\quad\Longleftrightarrow\quad
  \operatorname{supp}_A(x)\cap\operatorname{supp}_A(y)=\varnothing.
\]
This is Pitts' symmetric relation between elements of nominal sets over
the same atom carrier~\cite[Section~3.1, equation~(3.1)]{pitts2013}.
The value carriers may differ and have independent universes. Atom freshness
is a specialization: since $\operatorname{supp}_A(a)=\{a\}$,
\[
  a\mathrel{\#}x\quad\Longleftrightarrow\quad
  a\notin\operatorname{supp}_A(x).
\]

In Lean, \texttt{Fresh A x y} uses \texttt{Disjoint} on the two least supports,
avoiding a decidable-equality premise for forming an explicit intersection.
For individual certificates \texttt{hx} and \texttt{hy}, the corresponding
relation is \texttt{hx.FreshWith hy}. The two interfaces agree for any
certificates of the same elements under the same selected actions. The
atom helper \texttt{hx.Fresh a} specializes this relation using the canonical
atom certificate; \path{fresh_iff} retains the atom-to-carrier agreement law.
No additional support choice is made. Supported values in a non-nominal
carrier, such as pointwise constant functions, retain this freshness
interface without a nominality assertion about the whole function space.
These are named predicates; the mathematical symbol above is not an
additional Lean notation.

\begin{proposition}[General freshness laws]
\label{prop:general-freshness}
For finitely supported elements of the indicated selected actions,
\begin{align*}
 x\mathrel{\#}y&\quad\Longleftrightarrow\quad y\mathrel{\#}x,\\
 x\mathrel{\#}x&\quad\Longleftrightarrow\quad
   \operatorname{supp}_A(x)=\varnothing,\\
 (x,y)\mathrel{\#}z&\quad\Longleftrightarrow\quad
   x\mathrel{\#}z\ \land\ y\mathrel{\#}z,\\
 x\mathrel{\#}(y,z)&\quad\Longleftrightarrow\quad
   x\mathrel{\#}y\ \land\ x\mathrel{\#}z,\\
 (\pi\act x)\mathrel{\#}(\pi\act y)&\quad\Longleftrightarrow\quad x\mathrel{\#}y.
\end{align*}
Moreover, $x\mathrel{\#}y$ holds exactly when there exist disjoint finite
sets $S,T$ supporting $x,y$, respectively. Equivariant maps preserve
freshness in either argument: if $x\mathrel{\#}y$ and $f$ is equivariant,
then $f(x)\mathrel{\#}y$; likewise for a map on the second argument.
\end{proposition}

\begin{proof}
Symmetry and the self case are properties of disjointness. Exact product
support turns the product claims into disjointness of a union, while
least-support transport and injectivity of permutation images give the
renaming equivalence. The least supports themselves witness the finite-bound
characterization. Conversely, leastness puts the two supports inside any
disjoint sufficient bounds, and disjointness is preserved by taking subsets.
Finally, an equivariant image has support contained in that of its source,
so the same subset argument applies.
\end{proof}

The image implication generally does not reflect freshness: an equivariant
constant map into discrete data sends an atom $a$ to a value fresh for $a$,
although $a$ is not fresh for itself. The finite-bound characterization lets
clients establish freshness using convenient bounds without computing least
supports. These results reuse Mathlib's disjointness laws and the established
support formulas; their elementwise interfaces assume no nominality of the
whole carriers and no decidable atom equality. Classical equality needed for
finite images and unions remains inside proofs.

Products may now occur on both sides. For example,
$((a,b),c)\mathrel{\#}(d,e)$ decomposes into the six pairwise conditions
between $a,b,c$ and $d,e$. This follows by repeatedly applying the two
product equivalences, without any freshness-specific tactic.

\begin{proposition}[Freshness laws]
\label{prop:freshness-laws}
Let $x$ be finitely supported and $\pi\in\Perm(A)$.
If a finite set $S$ supports $x$, then $a\notin S$ implies
$a\mathrel{\#}x$. Moreover,
\begin{align*}
  \pi(a)\mathrel{\#}(\pi\act x)
    &\quad\Longleftrightarrow\quad a\mathrel{\#}x,\\
  a\mathrel{\#}(\pi\act x)
    &\quad\Longleftrightarrow\quad \pi^{-1}(a)\mathrel{\#}x.
\end{align*}
If $a\mathrel{\#}x$ and $b\mathrel{\#}x$, then
$\swap a b\act x=x$, including when $a=b$.
\end{proposition}

\begin{proof}
Leastness gives $\operatorname{supp}_A(x)\subseteq S$, so nonmembership
in $S$ suffices. Least-support transport gives
$\operatorname{supp}_A(\pi\act x)=\pi[\operatorname{supp}_A(x)]$.
Bijectivity of $\pi$ gives both membership equivalences and hence their
negations. Finally, both fresh endpoints avoid the supporting bound
$\operatorname{supp}_A(x)$, so the support swap criterion applies.
It already permits equal endpoints.
\end{proof}

The atom transport laws specialize the general freshness transport theorem.
Its formal proof combines the established support equation with Mathlib's
\path{Finset.disjoint_image} for injective images. It does not repeat the
conjugation argument. Although the support equation uses a finite-set
image, the freshness equivalences contain no such operation: classical
equality can stay inside their proofs. Their public statements therefore
need no decidable equality. The swap theorem retains the equality premise
of the swap operation, without adding a distinctness hypothesis.

Proposition~\ref{prop:canonical-support} yields the computation laws
\begin{gather*}
  a\mathrel{\#}b\ \Longleftrightarrow\ a\ne b,\qquad
  a\mathrel{\#}\operatorname{mk}(d),\qquad
  a\mathrel{\#}S\ \Longleftrightarrow\ a\notin S,\\
  a\mathrel{\#}(x,y)\ \Longleftrightarrow\
    a\mathrel{\#}x\ \land\ a\mathrel{\#}y.
\end{gather*}
The product law applies to individually supported components. Its Lean
statement likewise avoids an equality premise by keeping the union inside
the proof. Repeated decomposition handles either association of nested
products, including repeated components. For finite-set contexts, the
selected image action still requires decidable equality and the
\texttt{Pointwise} scope.
More generally, for a finite atom set $S$,
\[
 S\mathrel{\#}x\quad\Longleftrightarrow\quad
 (\forall a\in S,\ a\mathrel{\#}x)
 \quad\Longleftrightarrow\quad x\mathrel{\#}S.
\]
Since $\operatorname{supp}_A(S)=S$, this is disjointness expressed through
membership. Finite atom sets thus serve as contexts on either side of the relation.

\begin{proposition}[Fresh existence]
\label{prop:fresh-existence}
For every finitely supported $x$ and finite exclusion set $S\subseteq A$,
\[
  \exists a\in A,\qquad a\notin S\ \land\ a\mathrel{\#}x.
\]
In particular every such $x$ has a fresh atom. No nominality of its whole
carrier is required.
\end{proposition}

\begin{proof}
Choose an atom outside the finite set
$S\cup\operatorname{supp}_A(x)$. Infinite atoms guarantee its existence;
nonmembership in the two components gives the conclusion. Taking $S$ empty
gives the special case.
\end{proof}

Mathlib supplies this finite avoidance step as
\path{Finset.exists_notMem}. The elementwise Lean theorem
\path{FinitelySupported.exists_fresh_notMem} and its carrier adapter
\path{exists_fresh_notMem} keep the union inside the proof, so neither
requires decidable equality. Avoiding an explicit finite set alone uses
Mathlib's theorem directly and needs no action or nominality. The hypotheses
of a Finset action enter only when a finite set is treated as a renamed
context component.

For example, take
$c=(x,(y,(T,\operatorname{mk}(d))))$, where $x$ and $y$ are supported,
$T$ is a finite atom set and $d$ is arbitrary discrete data. Choose $a$
outside an explicit $S$ and fresh for $c$. Product decomposition yields
\[
  a\notin S,\qquad a\mathrel{\#}x,\qquad
  a\mathrel{\#}y,\qquad a\notin T.
\]
The same argument works using individual support certificates for the
components. There is no assumption that every element of their carriers
is supported. Alternatively, supplied sufficient bounds can be combined
inside the finite exclusion set and used via Proposition~\ref{prop:freshness-laws}.

Choose a second atom $b$ fresh for $c$ and outside $S\cup\{a\}$. Then
$a\ne b$, both avoid $S$ and $T$, and the fresh-swap law gives
$\swap a b\act c=c$, as well as fixation of $x$ and $y$ separately.
This combines fresh existence with an actual renaming consequence, without
requiring a new container operation or a freshness tactic.

These results assert existence and use witnesses within proofs. They supply
neither an executable fresh-name algorithm nor an equivariance or finite-support
property for a function choosing witnesses. The infinite-atom assumption
also remains essential: an explicit exclusion set may exhaust a finite
atom carrier, regardless of any nominality certificate.
